% Part: second-order-logic
% Chapter: sol-and-sets
% Section: introduction

\documentclass[../../../include/open-logic-section]{subfiles}

\begin{document}

\olfileid{sol}{set}{int}

\olsection{పరిచయం}

ద్వితీయ-స్థాయి తర్కంలో ప్రమేయాలతో పాటు వ్యక్తి క్షేత్రపు ఉపసమితులపై
కూడా పరిమాణీకరించవచ్చు కాబట్టి, సమితి సిద్ధాంతాన్ని కనీసం కొంతవరకైనా
ద్వితీయ-స్థాయి తర్కంలో నిర్వహించగలమని ఆశించవచ్చు. “నిర్వహించడం” అంటే,
సమితుల కోసం ప్రత్యేకమైన తార్కికేతర పదజాలం (ఉదాహరణకు, సమితి సిద్ధాంతపు
సభ్యత్వ \tetoken{విధేయం}{!!{predicate}}) ఏదీ లేకుండానే, సమితి-సిద్ధాంత ధర్మాలను,
ప్రకటనలను ద్వితీయ-స్థాయి తర్కంలో వ్యక్తం చేయగలమని మా ఉద్దేశం.
ఉదాహరణకు, సమితుల సమ్మేళనాలు, ఛేదనాలు, ఉపసమితి సంబంధాన్ని
నిర్వచించవచ్చు; అంతేకాక సమితుల పరిమాణాలను పోల్చవచ్చు, కాంటర్
సిద్ధాంతం వంటి ఫలితాలను పేర్కొనవచ్చు.

\end{document}
